% ---------------------------------------------------------------------------
% diagrams.tex — every TikZ figure in the deck, one macro each.
%
% Geometry note used throughout: a bearing is measured clockwise from north,
% while TikZ measures anticlockwise from east, so
%
%     tikz angle = 90 - bearing
%
% Bearing 038 -> 52 degrees, bearing 142 -> -52 degrees. The two beams
% therefore sit symmetrically either side of the east axis, which keeps every
% figure in this deck consistent with the survey's own orientation.
% ---------------------------------------------------------------------------

% === 1. The site: a triangular corner of a warehouse rooftop ===============
\newcommand{\diagRooftop}{%
\begin{tikzpicture}[scale=1]
  % roof slab + parapet
  \fill[figpale] (0,0) rectangle (9.6,5.6);
  \draw[figrule, line width=2pt] (0,0) rectangle (9.6,5.6);
  \draw[figrule, line width=0.6pt] (0.18,0.18) rectangle (9.42,5.42);

  % rooftop plant, to sell "this is a real roof, not a textbook"
  \foreach \x/\y/\w/\h in {6.4/3.7/1.5/1.2, 8.1/1.1/1.0/1.3, 6.6/1.2/0.9/0.7}{
    \fill[figmuted!25] (\x,\y) rectangle ++(\w,\h);
    \draw[figmuted!60, line width=0.6pt] (\x,\y) rectangle ++(\w,\h);
  }
  \draw[figmuted!45, line width=2.5pt] (7.15,3.7) -- (7.15,2.6) -- (8.6,2.6) -- (8.6,2.4);
  \node[dim, anchor=west] at (6.4,5.05) {existing HVAC plant};

  % north arrow
  \begin{scope}[shift={(0.85,4.35)}]
    \draw[figmuted, line width=0.9pt, -{Stealth[length=5pt]}] (0,0) -- (0,0.85);
    \node[dim, above] at (0,0.85) {N};
  \end{scope}

  % the survey triangle
  \coordinate (O) at (2.3,2.6);
  \coordinate (A) at ($(O)+(52:3.045)$);
  \coordinate (B) at ($(O)+(-52:2.394)$);
  \fill[figaccent!25] (O) -- (A) -- (B) -- cycle;
  \draw[figaccent!55, line width=0.8pt] (A) -- (B);
  \draw[beam] (A) -- (O) -- (B);

  \foreach \p in {O,A,B}{\node[vertex] at (\p) {};}
  \node[vlabel, left=3pt]        at (O) {O};
  \node[vlabel, above right=1pt] at (A) {A};
  \node[vlabel, below right=1pt] at (B) {B};

  % side labels ride the outside of each beam, clear of the vertex at O
  \node[dim, rotate=52,  anchor=south, yshift=1pt] at ($(O)!0.62!(A)$) {beam OA, 34.6 m};
  \node[dim, rotate=-52, anchor=north, yshift=-1pt] at ($(O)!0.62!(B)$) {beam OB, 27.2 m};
  \node[figink, font=\scriptsize\bfseries, align=center]
    at ($(O)!0.72!($(A)!0.5!(B)$)$) {unused\\corner};

  \node[dim, anchor=south west] at (0.25,-0.55) {Plan view of the warehouse rooftop, bounded by two structural beams};
\end{tikzpicture}}

% === 2. What a bearing actually is ========================================
\newcommand{\diagBearingConvention}{%
\begin{tikzpicture}[scale=1]
  \foreach \cx/\mth/\lab/\name in {0/52/038/A, 4.4/-52/142/B}{
    \begin{scope}[shift={(\cx,0)}]
      \draw[figrule, line width=0.8pt] (0,0) circle (1.45);
      \foreach \a/\t in {90/N, 0/E, -90/S, 180/W}{
        \draw[figrule, line width=0.8pt] (\a:1.32) -- (\a:1.45);
        \node[dim] at (\a:1.72) {\t};
      }
      \draw[north] (0,0) -- (90:1.32);
      \draw[ray, -{Stealth[length=5pt]}] (0,0) -- (\mth:1.32);
      \node[vertex] at (0,0) {};
      \node[vlabel, below left=1pt] at (0,0) {O};
      \node[vlabel] at (\mth:1.72) {\name};
      % clockwise sweep from north
      \draw[angmark, -{Stealth[length=4pt]}] (90:0.72) arc[start angle=90, end angle=\mth, radius=0.72];
      \node[figalert, font=\small\bfseries] at ({(90+\mth)/2}:1.05) {$\lab^\circ$};
    \end{scope}}
  \node[dim, align=center, text width=12cm] at (2.2,-2.35)
    {A bearing is always measured \emph{clockwise from north}, from the same origin each time.};
\end{tikzpicture}}

% === 3. Subtracting the bearings to get the included angle ================
\newcommand{\diagBearingSubtraction}{%
\begin{tikzpicture}[scale=1]
  \coordinate (O) at (0,0);
  \coordinate (A) at (52:3.9);
  \coordinate (B) at (-52:3.07);

  \draw[north] (O) -- (90:3.3);
  \node[dim, above] at (90:3.3) {N};

  \fill[figalert!14] (O) -- (52:1.05) arc[start angle=52, end angle=-52, radius=1.05] -- cycle;

  \draw[beam] (O) -- (A);
  \draw[beam] (O) -- (B);
  \node[vertex] at (O) {}; \node[vertex] at (A) {}; \node[vertex] at (B) {};
  \node[vlabel, below left=1pt]  at (O) {O};
  \node[vlabel, above right=1pt] at (A) {A};
  \node[vlabel, below right=1pt] at (B) {B};
  \node[dim, rotate=52,  anchor=south] at ($(O)!0.85!(A)$) {34.6 m};
  \node[dim, rotate=-52, anchor=north] at ($(O)!0.85!(B)$) {27.2 m};

  % bearing arcs, nested at different radii so they never collide
  \draw[ray, -{Stealth[length=4pt]}] (90:1.35) arc[start angle=90, end angle=52, radius=1.35];
  \node[figaccent, font=\small\bfseries] at (71:1.66) {$038^\circ$};
  \draw[ray, -{Stealth[length=4pt]}] (90:2.15) arc[start angle=90, end angle=-52, radius=2.15];
  \node[figaccent, font=\small\bfseries] at (14:2.48) {$142^\circ$};

  % the included angle: the thing we are actually after
  \draw[angmark] (52:1.05) arc[start angle=52, end angle=-52, radius=1.05];
  \node[figalert, font=\bfseries] at (0:0.72) {$104^\circ$};
\end{tikzpicture}}

% === 4. The general rule, including the case that needs adjusting =========
\newcommand{\diagBearingRule}{%
\begin{tikzpicture}[scale=0.95]
  % --- case 1: difference under 180, take it as it is
  \begin{scope}
    \draw[figrule, line width=0.8pt] (0,0) circle (1.5);
    \draw[north] (0,0) -- (90:1.5); \node[dim, above] at (90:1.5) {N};
    \fill[figalert!14] (0,0) -- (52:0.85) arc[start angle=52, end angle=-52, radius=0.85] -- cycle;
    \draw[ray, -{Stealth[length=4pt]}] (0,0) -- (52:1.5);
    \draw[ray, -{Stealth[length=4pt]}] (0,0) -- (-52:1.5);
    \draw[angmark] (52:0.85) arc[start angle=52, end angle=-52, radius=0.85];
    \node[figalert, font=\small\bfseries] at (0:1.22) {$104^\circ$};
    \node[dim] at (52:1.78) {$038^\circ$};
    \node[dim] at (-52:1.78) {$142^\circ$};
    \node[figink, font=\small, align=center] at (0,-2.35)
      {$142-38 = 104 < 180$\\[2pt] take it directly};
  \end{scope}
  % --- case 2: difference over 180, take the reflex complement
  \begin{scope}[shift={(6.4,0)}]
    \draw[figrule, line width=0.8pt] (0,0) circle (1.5);
    \draw[north] (0,0) -- (90:1.5); \node[dim, above] at (90:1.5) {N};
    % bearings 040 and 300 -> tikz 50 and -210 (=150)
    \fill[figalert!14] (0,0) -- (50:0.85) arc[start angle=50, end angle=150, radius=0.85] -- cycle;
    \draw[ray, -{Stealth[length=4pt]}] (0,0) -- (50:1.5);
    \draw[ray, -{Stealth[length=4pt]}] (0,0) -- (150:1.5);
    \draw[angmark] (50:0.85) arc[start angle=50, end angle=150, radius=0.85];
    \node[figalert, font=\small\bfseries] at (100:1.15) {$100^\circ$};
    \node[dim] at (50:1.8) {$040^\circ$};
    \node[dim] at (150:1.85) {$300^\circ$};
    \node[figink, font=\small, align=center] at (0,-2.35)
      {$300-40 = 260 > 180$\\[2pt] use $360 - 260 = 100^\circ$};
  \end{scope}
\end{tikzpicture}}

% === 5. The oblique triangle, reduced to SAS ==============================
\newcommand{\diagSAS}{%
\begin{tikzpicture}[scale=1]
  \coordinate (O) at (0,0);
  \coordinate (A) at (52:3.29);
  \coordinate (B) at (-52:2.58);

  \fill[figaccent!12] (O) -- (A) -- (B) -- cycle;
  \draw[beam] (O) -- (A);
  \draw[beam] (O) -- (B);
  \draw[figmuted, line width=1pt, dashed] (A) -- (B);

  \draw[angmark] (52:0.8) arc[start angle=52, end angle=-52, radius=0.8];
  \node[figalert, font=\bfseries] at (0:1.14) {$104^\circ$};
  \node[dim, anchor=east] at (0:0.42) {$C$};

  \foreach \p in {O,A,B}{\node[vertex] at (\p) {};}
  \node[vlabel, left=3pt]        at (O) {O};
  \node[vlabel, above right=1pt] at (A) {A};
  \node[vlabel, below right=1pt] at (B) {B};

  \node[figink, font=\small\bfseries, rotate=52,  anchor=south] at ($(O)!0.58!(A)$) {$b = 34.6$ m};
  \node[figink, font=\small\bfseries, rotate=-52, anchor=north] at ($(O)!0.58!(B)$) {$a = 27.2$ m};
  \node[dim, anchor=west, align=left, text width=2.6cm] at ($(A)!0.5!(B)+(0.3,0)$)
    {side AB, never measured and never needed};
\end{tikzpicture}}

% === 6. Where sin C comes from: it is the height you never measured =======
\newcommand{\diagHeight}{%
\begin{tikzpicture}[scale=1]
  \coordinate (O) at (0,0);
  \coordinate (B) at (2.86,0);
  \coordinate (A) at (104:3.63);
  \coordinate (F) at (-0.879,0);

  \fill[figaccent!12] (O) -- (A) -- (B) -- cycle;
  \draw[helper, dashed] (F) -- (O);          % base line, extended past O
  \draw[beam] (A) -- (O) -- (B);
  \draw[figmuted, line width=1pt, dashed] (A) -- (B);

  % the perpendicular height
  \draw[figaccent, line width=1.2pt, dashed] (A) -- (F);
  \draw[figaccent, line width=0.8pt] ($(F)+(0.22,0)$) -- ($(F)+(0.22,0.22)$) -- ($(F)+(0,0.22)$);
  \node[figaccent, font=\small\bfseries, rotate=90, anchor=south] at ($(A)!0.5!(F)$)
    {$h = b\sin C$};

  % angles at O
  \draw[angmark] (0:0.8) arc[start angle=0, end angle=104, radius=0.8];
  \node[figalert, font=\small\bfseries] at (52:1.12) {$104^\circ$};
  \draw[figaccent, line width=1pt] (104:0.42) arc[start angle=104, end angle=180, radius=0.42];
  \node[figaccent, font=\scriptsize\bfseries] at (142:0.68) {$76^\circ$};

  \foreach \p in {O,A,B}{\node[vertex] at (\p) {};}
  \node[vlabel, below right=1pt] at (O) {O};
  \node[vlabel, above right=1pt] at (A) {A};
  \node[vlabel, below right=1pt] at (B) {B};
  \node[dim, below=3pt] at ($(O)!0.55!(B)$) {$a = 27.2$ m};
  % right-aligned so the label runs away from O rather than over it
  \node[dim, anchor=north east] at ($(F)+(0.18,-0.14)$) {foot of the height};

  % Two separate nodes rather than one: TikZ's align=left redefines \\, and a
  % \\[dim] nested inside a {...} group within the node text breaks it.
  \node[figink, anchor=north west, font=\small, align=left] at (3.5,3.35)
    {Area $=\tfrac12 \times$ base $\times$ height\\[3pt]
     $\phantom{\text{Area}}=\tfrac12\, a\,(b\sin C)$};
  \node[figmuted, anchor=north west, font=\scriptsize, align=left] at (3.5,2.3)
    {$C$ is obtuse, so the foot lands\\[1pt]
     outside the triangle, and\\[1pt]
     $\sin 104^\circ = \sin 76^\circ$ still holds.};
\end{tikzpicture}}

% === 7. Illustrative panel layout (schematic, not an engineering layout) ==
\newcommand{\diagPanels}{%
\begin{tikzpicture}[scale=1]
  \coordinate (O) at (0,0);
  \coordinate (A) at (52:3.29);
  \coordinate (B) at (-52:2.58);

  \fill[figpale] (O) -- (A) -- (B) -- cycle;
  \begin{scope}
    \clip (O) -- (A) -- (B) -- cycle;
    \foreach \i in {0,...,8}{
      \foreach \j in {0,...,21}{
        \fill[panelcell] ({0.12+\i*0.26},{-2.3+\j*0.25}) rectangle ++(0.21,0.19);
      }}
  \end{scope}
  \draw[beam] (A) -- (O) -- (B);
  \draw[figmuted, line width=1pt] (A) -- (B);
  \foreach \p in {O,A,B}{\node[vertex] at (\p) {};}
  \node[vlabel, left=2pt]        at (O) {O};
  \node[vlabel, above right=1pt] at (A) {A};
  \node[vlabel, below right=1pt] at (B) {B};
  \node[dim, align=center, anchor=north] at (1.6,-3.0) {schematic only, ignoring tilt, row\\spacing and shading};
\end{tikzpicture}}
